% ======================================================================
%  Section 4.4 and 4.4.1 --- Principle.  Corrected.
%
%  TWO CLAIMS THAT ARE NO LONGER TRUE, AND ONE THAT IS NOW PROVABLE AND
%  MISSING.
%
%  A. "obtained from response maps that the scheme computes in any case,
%     the only additional work being a pointwise energy evaluation and a
%     mean over the image." False after 4.4.2. The feature vector uses the
%     SUMMED kernels; the structure measure uses the FIRST-ORDER kernels,
%     which the implementation evaluates in six additional convolutions per
%     iteration. This is the last of the three places the claim appears.
%
%  B. "It is not a generic image-gradient or saliency measure." With
%     first-order kernels it is precisely a Gauss-derivative gradient
%     magnitude, which is the canonical generic image-gradient measure.
%     4.4.2 already states the narrower claim that survives; 4.4.1 must
%     agree with it.
%
%  C. MISSING, AND IT IS THE BEST ARGUMENT THE SECTION COULD MAKE.
%     The sentence "where the response energy is high, the measurement
%     constrains the camera pose strongly" is asserted as an intuition. It
%     is in fact a theorem about the interaction matrix, and one that holds
%     for the first-order measure and NOT for the summed one. Every
%     component of the luminance interaction matrix row is a linear
%     combination of I_x and I_y with coefficients depending only on the
%     image position:
%
%       vvx = I_x/Z                     vwx = -xy I_x - (1+y^2) I_y
%       vvy = I_y/Z                     vwy = (1+x^2) I_x + xy I_y
%       vvz = -(x I_x + y I_y)/Z        vwz = y I_x - x I_y
%
%     so the norm of that row factorises as ||grad I|| times a purely
%     geometric function of (x,y). G is a scale-aggregated ||grad I||.
%     It therefore captures exactly the image-dependent factor of the row
%     norm, the remaining factor being known in closed form. That is why the
%     measure is the right one, and why the DC-carrying version was not: a
%     local mean of I bears no relation to the row norm at all.
%
%     Adding this converts the section's central intuition into an
%     identity, and retrospectively justifies the kernel choice of 4.4.2 on
%     grounds other than "the maps look better".
%
%  D. The band-pass sentence is now TRUE of the implementation and should be
%     kept. It was false while the summed kernels were in use.
% ======================================================================

\section{Hermite-Informed Weighting}
\label{sec:hermite-weight}
%%% The label belongs immediately after the \section command. In the source
%%% supplied it sits before \subsection{Principle} with no \section line
%%% above it, which would attach it to whatever precedes.

\subsection{Principle}
\label{sec:hermite-principle}

Remark~\ref{rem:weight-information} observed that the weights entering the
robust control law need not be functions of the residuals alone. The proposal
of this chapter is to make them depend also on the local image structure, as
measured in the Hermite representation on which the servoing already operates.

The motivation is the locality property established in
Section~\ref{sec:hermite-properties}. Each Hermite response depends on the
image only within a neighbourhood of radius of the order of the analysis scale
$\sigma$, so the response at a given position is a statement about the image
structure in that neighbourhood and at that scale.

%%% (C) The new paragraph. It replaces the first of the two "facts" with the
%%% reason that fact is true.
That the magnitude of the local response measures how strongly a measurement
constrains the pose is not merely plausible; it follows from the form of the
interaction matrix. Every component of the row of $\mathbf{L}_{R}$ associated
with a given position is a linear combination of the two image derivatives
$I_{x}$ and $I_{y}$ at that position, with coefficients that depend on the
image coordinates and the depth but not on the image content. The norm of that
row therefore factorises into the local gradient magnitude and a purely
geometric term known in closed form. The structure measure defined in the
following subsection is a scale-aggregated gradient magnitude, and so captures
exactly the image-dependent factor of the row norm --- precisely the part that
the residual alone cannot supply. Where it is large the measurement carries
strong pose information; where it is small the corresponding rows of
$\mathbf{L}_{R}$ are small, and the measurement constrains the motion weakly
whatever its residual happens to be.

%%% (A) The cost, stated honestly. Six convolutions is not "no additional
%%% work", and a reader who opens the implementation will count them.
The second consideration is cost. The measure is not free: the feature vector
is built from the summed Hermite kernels, whereas the structure measure uses
the first-order kernels of the same family, which the scheme must evaluate
separately. Three scales and two orientations add six convolutions per
iteration to the twenty-eight the feature and interaction maps already
require, together with a pointwise magnitude and a mean over the image. The
overhead is therefore of the order of a fifth of the filtering already
performed, and no new representation, scale or parameter is introduced.

%%% (B) The claim narrowed to what survives, matching the wording adopted in
%%% 4.4.2. (D) The band-pass clause kept: it is now true of the code.
The essential point is that this reliability assessment is performed \emph{in
the same representation in which the control operates}. The operator itself is
a classical one --- a Gaussian-derivative gradient magnitude --- and no claim
is made for its novelty. What distinguishes it is that it is evaluated at the
same analysis scales that define the features, and is drawn from the same
Gauss--Hermite family, so that it inherits the band-pass character and the
noise suppression established in Section~\ref{sec:hermite-properties}: the
structure it reports is structure at the scales the controller uses, with the
high-frequency sensor noise already attenuated by the Gaussian envelope of the
filters.

% ======================================================================
%  WHERE THE THREE "ALREADY COMPUTED" CLAIMS NOW STAND
%    4.1  introduction   -- corrected, states the six convolutions
%    4.3.3               -- corrected, claims only that the representation
%                           is already in use
%    4.4.1  here         -- corrected, states the operation count
%  If you keep only one of the three quantitative statements, keep this one;
%  the other two can simply avoid the claim.
%
%  CONSISTENCY CHECK AFTER THIS EDIT
%  4.4.2 says the measure is "obtained from kernels the scheme already
%  evaluates". That phrase is from my own earlier draft and has the same
%  defect. Replace it there with the wording used above, or delete the
%  clause: the sentence stands without it.
% ======================================================================
